% Part: computability
% Chapter: computability-theory
% Section: universal-part-function

\documentclass[../../../include/open-logic-section]{subfiles}

\begin{document}

\olfileid{cmp}{thy}{uni}
\olsection{সার্বজনীন আংশিক গণনসাধ্য অপেক্ষক}

\begin{thm}
\ollabel{thm:univ-comp}
একটি সার্বজনীন আংশিক গণনসাধ্য অপেক্ষক~$\fn{Un}(e,x)$ আছে। অন্যভাবে
বললে, এমন একটি অপেক্ষক $\fn{Un}(e,x)$ আছে যাতে:
\begin{enumerate}
\item $\fn{Un}(e,x)$ আংশিক গণনসাধ্য।
\item $f(x)$ যেকোনো আংশিক গণনসাধ্য অপেক্ষক হলে এমন একটি স্বাভাবিক
সংখ্যা $e$ আছে, যাতে প্রত্যেক~$x$-এর জন্য
$f(x) \simeq \fn{Un}(e,x)$।
\end{enumerate}
\end{thm}

\begin{proof}
$\fn{Un}(e,x) \simeq U(\umin{s}{T(e,x,s)})$ ধরা যাক, যেখানে $U$ ও $T$
ক্লিনির স্বাভাবিক রূপের উপপাদ্যে
(\olref[nfm]{thm:normal-form}) দেওয়া অপেক্ষক ও সম্বন্ধ।
\end{proof}

\begin{explain}
এটি শুধু এই কথার নির্ভুল রূপ যে আংশিক গণনসাধ্য অপেক্ষকগুলির একটি কার্যকর
তালিকায়ন আমাদের আছে। ধারণাটি হলো, $f_e(x) = \fn{Un}(e,x)$ সমীকরণে
সংজ্ঞায়িত অপেক্ষককে $f_e$ লিখলে $f_0$, $f_1$, $f_2$, \dots অনুক্রমে
সব আংশিক গণনসাধ্য অপেক্ষক থাকে এবং~$e$ ও~$x$ থেকে “একই নিয়মে”
$f_e(x)$ গণনা করা যায়। সরলতার জন্য আমরা এখানে একস্থানী অপেক্ষকের জন্য
সার্বজনীন একটি দ্বিস্থানী অপেক্ষক ব্যবহার করছি; কিন্তু সংখ্যার অনুক্রম
সংকেতায়িত করে একে আরও বেশি আর্গুমেন্টের ক্ষেত্রে সহজেই সাধারণীকরণ করা
যায়। যেমন, $f(x,y,z)$ কোনো $3$-স্থানী আংশিক পুনরাবৃত্ত অপেক্ষক হলে
$g(x) \simeq f((x)_0, (x)_1, (x)_2)$ একটি একস্থানী আংশিক পুনরাবৃত্ত
অপেক্ষক।
% BN-SRC-143 TARGET-CORRECTION-BEGIN
\par\smallskip\noindent\textnormal{\small\textbf{উৎস-সংশোধন (BN-SRC-143):} হিমায়িত ইংরেজি পাঠের শেষ বাক্যে $g$-কে “unary recursive function” বলা হয়েছে, যদিও অনুমিত~$f$ আংশিক এবং সংজ্ঞাটি~$g$-এর সর্বত্র-সংজ্ঞায়িততা নিশ্চিত করে না। বাংলা পাঠে সঠিকভাবে “একস্থানী আংশিক পুনরাবৃত্ত অপেক্ষক” লেখা হয়েছে।} % BN-SRC-143 TARGET-CORRECTION-END
\end{explain}

\end{document}
